\subsection{Consequences for a Hamiltonian threshold theory}
\label{threshold:negative-conclusions}

The obstruction side supports the following precise interpretation.
Temperature changes the weight and the noise level of latent sectors;
full AA fails when a sector retains a quantum-accessible hidden label
whose classical identification or decision cost grows with system
size, at an accuracy where the sector cannot be ignored.  The
two-level code gives an energy--entropy threshold.  The fixed soft-code
Hamiltonian gives a contrasting endpoint: its global code order
survives every fixed positive inverse temperature, despite arbitrarily
weak one-spin energy increments.

Several tempting structural tests are therefore inadequate even for
the full query target, not just for influence AA:
\begin{itemize}
\item Identity covariance, uniform three-spin marginals, and entropy
      density arbitrarily close to $\log2$ do not suffice.
\item A fixed bound on every single-spin energy change and an $O(n)$
      energy range do not suffice.
\item Merely observing $w_n\to0$ for a hard Gibbs sector does not
      suffice without relating this decay to the classical query cost.
\item Phase coexistence or a low-degree low-influence observable does
      not by itself establish a full-AA failure.
\end{itemize}

None of these examples establishes failure for fixed-arity,
bounded-incident-strength interactions satisfying a specified
high-temperature contraction criterion.  In particular a bound on
$|H(x)-H(x^{(i)})|$ is weaker than a bound on a Dobrushin influence
row sum: a small range of a conditional bias does not bound the sum
of its separate responses to many other coordinates.  Nor does the
soft-code formula represent a bounded-degree interaction graph.
Those remain meaningful candidate hypotheses for positive full-AA
theorems.

The quantum primitive is taken from
\href{https://people.eecs.berkeley.edu/~vazirani/pubs/bv.pdf}
{Bernstein--Vazirani, \emph{Quantum complexity theory} (1997)}.
The broader relevance of coherent Fourier structure to query
separations is demonstrated by
\href{https://arxiv.org/abs/1411.5729}
{Aaronson--Ambainis, \emph{Forrelation} (2014 preprint; 2018 journal)};
no Forrelation lower bound is needed in the proofs here.  The
positive comparisons import uniform-input AA in precisely the sense
already stipulated by the shared notes.  The noisy separation, the
fixed-Hamiltonian comparison, and the arbitrary-weight criterion are
self-contained deductions above.
